\subsection{A degree-two class surviving the based PD saturation}

\begin{proposition}[A-L12: Proved; the saturated comparison can fail in degree two]
\label{AL12}
Use the resolution and the actual nonzero quadratic
example at the end of A-L11, with $b\geq3$ and $s=2b$.
Apply the iterated based PD saturation of A-SAT01
to its resolved-tail algebra, retaining its map to
the unresolved tail algebra.
The saturated resolved algebra has a class of exact
order two in degree two and weight $2s+5$, in
simplex two.  The actual comparison sends this
class to zero, and the unresolved tail algebra
has zero degree-two cohomology in that simplex.
Thus the saturated comparison need not be a
cohomology isomorphism through degree three.
This statement concerns the displayed comparison;
further degree-two repairs are not excluded.
\end{proposition}

\begin{proof}
We first work in the pure zero-twist algebra $E$
over $S=\Gamma_{\mathbb Z_{(2)}}(T)$ from A-L11.
Keep its bottom generators $c,h,k$, with
$Tc=2(k-h)$, and use the following names for
the free positive-form generators of $W$:
\[
 \begin{array}{c|c|c}
 \text{symbol}&\text{summand of }P&
                  \text{form degree, unshifted weight}\\\hline
 x&v_{r_2}&1,2\\
 y&v_{r_3}&1,3\\
 g&u_{g_3}&1,3\\
 r&u_{r_2}&2,2\\
 a&v_{e_1}&2,1\\
 b'&v_{e_2}&2,2
 \end{array}
\]
Every actual weight includes the common shift $s$.
The exact linear differential gives
\[
 dc=2g+Tx-2y,\qquad
 dx=2r+Ta-2b'.                                      \tag{AL.68}
\]
Define the integral two-form
\[
 \xi=gx-yx+cr-cb'+(k-h)a.                           \tag{AL.69}
\]
All terms have arity two and weight $2s+5$.
Using alternating odd squares and (AL.61), one obtains
\[
 \begin{split}
 d(cx)&=(2g+Tx-2y)x+c(2r+Ta-2b')\\
      &=2\bigl(gx-yx+cr-cb'+(k-h)a\bigr)
        =2\xi.
 \end{split}                                       \tag{AL.70}
\]
Torsion-freeness and $d^2=0$ imply $d\xi=0$.
The $gx$ component is nonzero in the rational
exterior algebra on the independent free
form-one generators $g,x,y$.
Hence $\xi\ne0$.

The entire arity-two degree-one lattice in this
weight is $R\,cx$: a zero-form generator has
weight at least $s+3$, and a form-one generator
has weight at least $s+2$, with exactly one
generator at each minimum.  There is no remaining
coefficient weight.  All other choices have larger
weight.  Thus an integral primitive of $\xi$
would have to be $\alpha cx$ with
$\alpha\in R=\mathbb Z_{(2)}$.
Equation (AL.70) and nonvanishing give $2\alpha=1$,
which is impossible in $R$.
The arity grading prevents a primitive in another
arity from changing this conclusion.
Consequently $[\xi]$ has exact order two in $H^2(E)$.

We next check the entire based PD saturation,
including its products, rather than only its
first linear step.  At this fixed simplex define
the algebraic closure $E^{\mathrm{bs}}$ inside
$E_{\mathbb Q}$ by starting with $E$, adjoining
every rational zero-form whose differential is
already present and whose augmentation to
$S_{\mathbb Q}$ is integral, taking the generated
PD differential subalgebra, and repeating countably
many times.  This is the pointwise operation in
A-SAT01 with the augmentation in place of the
vertex-zero map.

Every stage is arity graded.  Indeed its derivative
condition separates by arity, and the augmentation
vanishes on every positive-arity summand.
Its arity-zero algebra stays $S$: a rational
coefficient has integral augmentation precisely
when it belongs to $S$.
The arity-one cochain complex is $W$, whose
degree-one cohomology is zero and whose degree-zero
differential is injective.  A new rational arity-one
zero-form with integral differential therefore
already belongs to $W^0$.  Products or divided
powers of elements of larger arity cannot change
this assertion at a later stage.
Thus arity one remains unchanged throughout.

All rational arity-two zero-forms have weight at
least $2(s+3)=2s+6$, since they are products or
divided squares of the zero-form generators of
minimum weight $s+3$.
Therefore no stage can add an arity-two zero-form
of weight at most $2s+5$.
Multiplication and divided powers cannot decrease
weight or arity.  Since arity zero and arity one
were unchanged, the entire arity-two cochain
complex in weight $2s+5$ is unchanged by every
stage and by their union.
In particular (AL.69) remains an exact-order-two
class in $H^2(E^{\mathrm{bs}})$.

Return now to the actual independent quadratic
example of A-L11.  In its simplex-two algebra,
the expressions (AL.68) still hold: the twist is
zero on $W^0$ and $W^1$.
Thus (AL.69)--(AL.70) give the same closed
two-form before saturation, even with the nonzero
quadratic twist on the top generators.
There is an actual augmented PD cochain map
from this simplex algebra to $E$: evaluate the
$G_R$ factor at vertex zero, kill the old
identity-disk variables, and retain the $W$
generators.  Killing those old variables kills
the twist, so the differential is respected.
The composite with $E\to S$ is exactly the
source vertex-zero map.  This is a map of the
specified simplex algebras; no simplicial
retraction $G_R\to S$ is claimed.

This map extends at every saturation stage into
the algebraic closure $E^{\mathrm{bs}}$.
To verify this directly, a newly adjoined rational
zero-form has derivative in the preceding source
stage; its image has derivative in the preceding
target closure stage by induction.  Its integral
vertex-zero value becomes its integral augmentation.
Its image is therefore one of the permitted new
zero-forms of the target closure.  Products and
divided powers then extend the map to the generated
subalgebra.  Induction and union prove the asserted
map after the full iterated saturation.
It sends (AL.69) to the nonzero class just proved.
Thus the actual saturated resolved algebra has
the same exact-order-two class.

Finally the augmentation comparison (AL.67) kills
$c,h,k,x,y,g,r$ and therefore kills (AL.69).
It extends through saturation by A-SAT01, including
the unresolved tail algebra among the specified
targets.  In simplex two that target is exactly
the torsion-free tail algebra of A-S09-21 for
the independent class $[\gamma_2(y_*)]$.
Its $H^2$ is zero by that proved calculation.
These are the actual maps and the resulting
failed degree-two comparison claimed in the
statement.
\end{proof}
